%==============================================================================
% FICHIER UNIQUE pour Overleaf — ouvrir UNIQUEMENT ce fichier comme document principal
% Menu Overleaf : Menu > Document principal > StoryTranslator_Rapport_Overleaf.tex
% Compilateur : PdfLaTeX
%==============================================================================
\documentclass[12pt,twoside,letterpaper]{article}

\usepackage[utf8]{inputenc}
\usepackage[T1]{fontenc}
\usepackage[french]{babel}
\usepackage{lmodern}
\usepackage{microtype}
\usepackage{geometry}
\geometry{margin=2.5cm}
\usepackage{setspace}
\onehalfspacing
\usepackage{graphicx}
\usepackage{xcolor}
\usepackage{hyperref}
\hypersetup{colorlinks=true,linkcolor=black,citecolor=blue,urlcolor=blue}
\usepackage{float}
\usepackage{enumitem}
\usepackage{amsmath,amssymb}
\usepackage{csquotes}

\newcommand{\reporttitle}{{\huge \textsc{Rapport} \\ \vspace{0.5cm} {\Large Projet de fin d'\'etudes --- StoryTranslator}}}
\newcommand{\reportauthorOne}{Aissatou Seck}
\newcommand{\reportauthorTwo}{}
\newcommand{\reportauthorThree}{}
\newcommand{\reportProfessor}{Yacine Yaddaden, Ph. D.}
\newcommand{\reporttype}{Coursework}

\begin{document}

% Last modification: 2016-09-29 (Marc Deisenroth)
% Modification for UW: 2017-05-22 (jphickey)
\begin{titlepage}

\newcommand{\HRule}{\rule{\linewidth}{0.5mm}} % Defines a new command for the horizontal lines, change thickness here


%----------------------------------------------------------------------------------------
%	LOGO SECTION
%----------------------------------------------------------------------------------------



\begin{center} % Center remainder of the page

%----------------------------------------------------------------------------------------
%	HEADING SECTIONS
%----------------------------------------------------------------------------------------

\includegraphics[width=7cm]{uqar}\\[1.5cm]
\textbf{\textsc{\Large PROJET DE FIN D'ETUDES EN INFORMAT\\ \vspace{0.5cm} INF39615-7T}}\\[1.0cm]
\textsc{\Large Universit\'e du Qu\'ebec \`a Rimouski}\\[0.5cm]
\textsc{\large D\'epartement de math\'ematiques, d'informatique et de g\'enie}\\[0.95cm]

%----------------------------------------------------------------------------------------
%	TITLE SECTION
%----------------------------------------------------------------------------------------

\HRule \\ % [0.4cm]
{ \huge \bfseries \reporttitle}\\ % Title of your document
\HRule \\[1.5cm]
\end{center}
%----------------------------------------------------------------------------------------
%	AUTHOR SECTION
%----------------------------------------------------------------------------------------

%\begin{minipage}{0.4\hsize}
\begin{flushleft} \large
\textit{\textsc{\textbf{\'Equipe :}}}\\
\reportauthorOne \\ %~(Matricule: \cidOne)\\ %
% \reportauthorTwo \\
% \reportauthorThree \\
\vspace{0.5cm}
\textit{\textsc{\textbf{Professeur :}}}\\
\reportProfessor
\end{flushleft}
% \vspace{3.5cm}
\vfill
\makeatletter
Date: \@date

% \vfill % Fill the rest of the page with whitespace



\makeatother


\end{titlepage}

\tableofcontents
\newpage

%========================================================
\section{Introduction}
%========================================================

Ce rapport pr\'esente la poursuite du projet \textbf{StoryTranslator},
une application web destin\'ee \`a aider les enfants \`a apprendre une langue
par la lecture d'histoires illustr\'ees.

Une premi\`ere version avait \'et\'e r\'ealis\'ee dans le cadre du cours
Projet en informatique II (INF34615-7T).
Elle permettait d\'ej\`a de lire des histoires en plusieurs langues,
d'\'ecouter le texte \`a voix haute et de sauvegarder la progression
dans le navigateur.

Dans le cadre du projet de fin d'\'etudes (INF39615-7T), l'application a \'et\'e enrichie
pour devenir une plateforme plus compl\`ete :
cr\'eation de comptes utilisateurs, sauvegarde des donn\'ees en base,
espace administrateur, interface selon la langue maternelle,
et quiz de compr\'ehension assist\'e par l'intelligence artificielle,
avec un mode oral fond\'e sur Whisper.

Le code source du projet est disponible \`a l'adresse suivante :
\url{https://github.com/Seck2000/StoryTranslatorWeb}.

Ce document d\'ecrit le fonctionnement du site, les technologies utilis\'ees,
puis les \'el\'ements techniques utiles \`a la maintenance et \`a l'\'evolution du projet.

\newpage

%========================================================
\section{Pr\'esentation du Projet}
%========================================================

\subsection{Authentification et comptes utilisateurs}

Avant d'acc\'eder \`a la biblioth\`eque, l'utilisateur doit cr\'eer un compte
ou se connecter.
\`A l'inscription, l'enfant (ou le parent) renseigne notamment :
pr\'enom, nom, courriel, langue maternelle, langue cible, niveau,
tranche d'\^age et mot de passe.

Les mots de passe sont hach\'es (bcrypt) et les sessions s'appuient sur des jetons JWT.
Deux r\^oles existent : \texttt{user} et \texttt{admin}.

\subsection{Page d'accueil : La Biblioth\`eque}

Apr\`es authentification, l'utilisateur arrive dans son espace personnel (espace enfant).
La biblioth\`eque affiche les histoires disponibles sous forme de cartes
(image de couverture, titre, progression).

Depuis cette page, l'utilisateur peut :
\begin{itemize}[leftmargin=*]
  \item cliquer sur une carte pour lancer la lecture ;
  \item consulter ses favoris et ses histoires r\'ecemment lues ;
  \item ouvrir son profil pour modifier ses informations ;
  \item lancer une histoire au hasard.
\end{itemize}

Les histoires affich\'ees d\'ependent de la \textbf{tranche d'\^age}
d\'eclar\'ee dans le profil (3 \`a 5 ans, 6 \`a 8 ans, 9 \`a 12 ans).

\subsection{Profil utilisateur et pr\'ef\'erences}

Dans \og Mon profil \fg{}, l'utilisateur peut modifier ses informations et sa photo.
La \textbf{langue maternelle} d\'etermine la langue des textes de l'interface.
La \textbf{langue cible} est celle que l'enfant apprend (quiz, r\'eponses attendues).

\subsection{Importation d'une Histoire}

L'importation d'histoires est r\'eserv\'ee \`a l'\textbf{administrateur}.
Celui-ci clique sur \og Importer \fg{} et s\'electionne un fichier \texttt{.zip} contenant :
\begin{itemize}[leftmargin=*]
  \item un fichier \texttt{story.json} (titre, sc\`enes, textes, chemins d'images) ;
  \item les images des sc\`enes (et \'eventuellement des avatars).
\end{itemize}

Le serveur d\'ecompresse le fichier, v\'erifie la pr\'esence de \texttt{story.json},
puis ajoute l'histoire \`a la biblioth\`eque.
Si le fichier est invalide, il est rejet\'e avec un message d'erreur clair.

\subsection{Espace administrateur}

L'administrateur dispose d'un tableau de bord distinct avec import ZIP bien visible.
Son profil est simplifi\'e : pr\'enom, nom, courriel et photo uniquement
(pas de langue maternelle, langue cible ni niveau).

\subsection{Le Lecteur d'Histoire}

Une fois une histoire s\'electionn\'ee, l'interface bascule en mode lecteur.
L'utilisateur voit :
\begin{itemize}[leftmargin=*]
  \item l'image de la sc\`ene ;
  \item l'avatar du personnage ;
  \item la bulle de texte dans la langue s\'electionn\'ee ;
  \item une barre / un compteur de progression ;
  \item les boutons de navigation (pr\'ec\'edent / suivant) ;
  \item sur la derni\`ere sc\`ene, un bouton pour discuter avec l'IA.
\end{itemize}

\subsubsection{Changement de Langue}
Un contr\^ole permet de changer la langue d'affichage du texte de l'histoire.
Lorsque l'arabe est s\'electionn\'e, le texte peut s'afficher de droite \`a gauche (RTL).

\subsubsection{Contr\^oles Audio}
Le site utilise la synth\`ese vocale du navigateur (\texttt{SpeechSynthesis})
pour lire le texte \`a voix haute :
\begin{itemize}[leftmargin=*]
  \item lecture automatique \`a chaque changement de sc\`ene (activable / d\'esactivable) ;
  \item lecture manuelle (lecture / pause / reprise) ;
  \item indication de la langue de lecture au navigateur
  (ex.\ \texttt{fr-FR}, \texttt{en-US}, \texttt{ar-SA}),
  afin qu'il choisisse une voix compatible avec le texte affich\'e.
\end{itemize}

\subsubsection{Mode Immersif (Masquer le Texte)}
Un bouton permet de masquer la bulle de texte afin de favoriser la compr\'ehension orale.

\subsubsection{Sauvegarde de Progression}
La progression est sauvegard\'ee en base de donn\'ees pour l'utilisateur connect\'e.
\`A la r\'eouverture d'une histoire, une bo\^ite de dialogue propose de reprendre ou de recommencer.

\subsection{Quiz de fin d'histoire avec l'IA}

\`A la fin d'une histoire, l'enfant peut discuter avec une IA (Gemini).
Le tuteur :
\begin{itemize}[leftmargin=*]
  \item pose des questions uniquement sur le contenu de l'histoire ;
  \item demande un court r\'esum\'e ;
  \item corrige gentiment les erreurs de langue ;
  \item communique exclusivement dans la langue cible ;
  \item refuse les chiffres (r\'eponses en toutes lettres).
\end{itemize}

\subsection{Mode oral et transcription Whisper}

Le quiz propose un mode \textbf{\'Ecrire} et un mode \textbf{Parler}.
En mode Parler :
\begin{enumerate}[leftmargin=*]
  \item l'IA lit sa question \`a voix haute ;
  \item l'enfant enregistre sa r\'eponse au micro ;
  \item l'audio est envoy\'e au serveur ;
  \item \textbf{Whisper} (OpenAI) transcrit la voix mot pour mot ;
  \item le texte est transmis \`a l'IA pour correction et suite du quiz.
\end{enumerate}

Cette fid\'elit\'e de transcription est essentielle :
l'IA doit voir les vraies erreurs pour pouvoir les corriger.

\subsection{Interface multilingue (langue maternelle)}

Les menus, boutons et messages suivent la langue maternelle du profil,
m\^eme apr\`es d\'econnexion / reconnexion.
Dans les listes de langues, chaque langue appara\^it dans sa forme native
(ex.\ English, Espagnol, Arabe) pour \^etre facilement identifiable.

\newpage

%========================================================
\section{Technologies Utilis\'ees}
%========================================================

\subsection{Front-end (Interface Client)}
\begin{itemize}[leftmargin=*]
  \item \textbf{React.js} : construction de l'interface.
  \item \textbf{Vite} : outil de d\'eveloppement et de compilation.
  \item \textbf{Tailwind CSS} : mise en forme responsive.
  \item \textbf{Axios} : communication HTTP avec l'API.
  \item \textbf{Lucide-react} : ic\^ones.
  \item \textbf{Web Speech API} : synth\`ese vocale.
  \item \textbf{MediaRecorder} : enregistrement audio pour Whisper.
\end{itemize}

\subsection{Back-end (Serveur)}
\begin{itemize}[leftmargin=*]
  \item \textbf{Node.js} et \textbf{Express} : API REST.
  \item \textbf{Multer} / \textbf{Unzipper} : import ZIP.
  \item \textbf{CORS} : autorisation des appels client $\leftrightarrow$ serveur.
  \item \textbf{JWT} et \textbf{bcrypt} : authentification et mots de passe.
\end{itemize}

\subsection{Base de donn\'ees et authentification}
\begin{itemize}[leftmargin=*]
  \item \textbf{PostgreSQL} : persistance des comptes, favoris, progression, historique et pr\'ef\'erences.
  \item \textbf{Prisma} : sch\'ema et migrations.
\end{itemize}

\subsection{Intelligence artificielle et reconnaissance vocale}
\begin{itemize}[leftmargin=*]
  \item \textbf{Gemini} : quiz conversationnel et correction linguistique.
  \item \textbf{Whisper (OpenAI)} : transcription orale litt\'erale.
\end{itemize}

\subsection{Qualit\'e et tests}
\begin{itemize}[leftmargin=*]
  \item \textbf{Vitest} : tests unitaires du projet.
  \item \textbf{GitHub Actions} : int\'egration continue (CI).
  \`A chaque \texttt{push} ou \texttt{pull request},
  les tests sont ex\'ecut\'es automatiquement pour d\'etecter les erreurs rapidement.
  Il n'y a pas de d\'eploiement automatique (pas de CD).
\end{itemize}

\newpage

%========================================================
\section{Description Technique pour la Maintenance}
%========================================================

Cette section est destin\'ee \`a tout d\'eveloppeur souhaitant comprendre,
modifier ou faire \'evoluer le projet.

\subsection{Structure des Fichiers}

\begin{verbatim}
StoryTranslatorWeb/
  client/
    src/
      App.jsx
      components/
      utils/
      constants/
  server/
    index.js
    routes/
    services/
    prisma/
    uploads/
  tests/
  .github/workflows/
\end{verbatim}

\subsection{Le Fichier Client : App.jsx}

\texttt{App.jsx} reste le point central de navigation entre les vues :
accueil public, authentification, biblioth\`eque (enfant ou admin), lecteur, quiz.

\subsection{Le Fichier Serveur : index.js}

Le serveur Express branche plusieurs routes prot\'eg\'ees, notamment :
\begin{itemize}[leftmargin=*]
  \item \texttt{/api/auth} : inscription, connexion, profil, avatar ;
  \item \texttt{/api/library} : favoris, historique, progression ;
  \item \texttt{/api/ai} : d\'ebut et suite du quiz ;
  \item \texttt{/api/speech} : transcription Whisper ;
  \item \texttt{/api/stories} et upload ZIP : biblioth\`eque d'histoires.
\end{itemize}

\subsection{Le Format d'une Histoire : story.json}

Le format initial est conserv\'e. Exemple simplifi\'e :

\begin{verbatim}
{
  "id": "lina-parc",
  "title": "Une journee au parc",
  "thumbnail": "images/scene1.png",
  "ageBand": "petits",
  "scenes": [
    {
      "id": 1,
      "image": "images/scene1.png",
      "character": {
        "name": "Lina",
        "avatar": "images/scene1.png"
      },
      "text": {
        "fr": "Bonjour ! Je m'appelle Lina.",
        "en": "Hello! My name is Lina.",
        "ar": "..."
      }
    }
  ]
}
\end{verbatim}

\subsection{Comment Lancer le Projet}

\textbf{Pr\'erequis :} Node.js et PostgreSQL configur\'es,
fichier \texttt{server/.env} renseign\'e.

\begin{enumerate}[leftmargin=*]
  \item Terminal 1 : \texttt{cd server}, puis \texttt{npm install}, puis \texttt{node index.js}.
  \item Terminal 2 : \texttt{cd client}, puis \texttt{npm install}, puis \texttt{npm run dev}.
\end{enumerate}

Ouvrir l'URL Vite (ex.\ \texttt{http://localhost:5173}).

\subsection{Comment Modifier le Site}

\begin{itemize}[leftmargin=*]
  \item \textbf{Design} : classes Tailwind dans les composants React.
  \item \textbf{Nouvelle langue d'interface} : \texttt{translations.js}.
  \item \textbf{Nouvelle histoire} : ZIP conforme \`a \texttt{story.json}, import admin.
  \item \textbf{Prompt du quiz} : \texttt{server/services/aiChat.js}.
  \item \textbf{Transcription} : \texttt{server/services/whisperTranscribe.js}.
\end{itemize}

\subsection{Configuration des variables d'environnement}

Le fichier \texttt{server/.env} contient notamment :
\begin{itemize}[leftmargin=*]
  \item \texttt{DATABASE\_URL}
  \item \texttt{JWT\_SECRET}
  \item \texttt{GEMINI\_API\_KEY}
  \item \texttt{OPENAI\_API\_KEY} (Whisper)
  \item \texttt{WHISPER\_MODEL}
\end{itemize}
Ces secrets ne doivent pas \^etre commit\'es dans Git.

\newpage

%========================================================
\section{Conclusion}
%========================================================

\subsection{Bilan du Projet}

L'application \textbf{StoryTranslator} est pass\'ee d'un simple outil de lecture multilingue
\`a une plateforme plus compl\`ete.
Elle int\`egre d\'esormais des comptes s\'ecuris\'es, une persistance des donn\'ees,
des r\^oles utilisateur et administrateur, une interface adapt\'ee \`a la langue maternelle,
ainsi qu'un quiz guid\'e par l'IA avec un mode oral fond\'e sur Whisper.
L'architecture client/serveur initiale a \'et\'e conserv\'ee et \'etendue de mani\`ere coh\'erente.

Le d\'ep\^ot GitHub du projet est accessible ici :
\url{https://github.com/Seck2000/StoryTranslatorWeb}.

\subsection{Ce que ce Projet m'a Appris}

Ce projet m'a permis de consolider :
\begin{itemize}[leftmargin=*]
  \item React et la gestion d'\'etat ;
  \item l'architecture client/serveur et les API REST prot\'eg\'ees ;
  \item PostgreSQL et la mod\'elisation des donn\'ees utilisateur ;
  \item l'int\'egration de Gemini et Whisper ;
  \item les tests unitaires et l'int\'egration continue (CI) avec GitHub Actions.
\end{itemize}

\subsection{Am\'eliorations Futures}

\begin{itemize}[leftmargin=*]
  \item d\'eploiement en ligne (frontend, backend, base h\'eberg\'ee, HTTPS) ;
  \item mise en cache \'eventuelle avec Redis pour acc\'el\'erer certaines lectures ;
  \item renforcement des tests d'int\'egration ;
  \item \'editeur d'histoires en ligne ;
  \item am\'elioration continue de l'exp\'erience mobile.
\end{itemize}

\newpage

\begin{thebibliography}{99}
\bibitem{whisper2022}
Alec Radford et al.
\newblock Robust Speech Recognition via Large-Scale Weak Supervision.
\newblock \emph{arXiv:2212.04356}, 2022.
\newblock \url{https://openai.com/research/whisper}.

\bibitem{gemini}
Google.
\newblock Gemini API.
\newblock \url{https://ai.google.dev/}, 2024.

\bibitem{react}
Meta.
\newblock React Documentation.
\newblock \url{https://react.dev/}, 2024.

\bibitem{express}
OpenJS Foundation.
\newblock Express Documentation.
\newblock \url{https://expressjs.com/}, 2024.
\end{thebibliography}

\end{document}
